\section{Challenges and Future Directions}
\label{sec:future}


The development of MLLMs is still in a rudimentary stage and thus leaves much room for improvement, which we summarize below:


\begin{itemize}

\item Current MLLMs are limited in processing multimodal information of long context. This restricts the development of advanced models with more multimodal tokens, \eg long-video understanding, and long documents interleaved with images and text.

\item MLLMs should be upgraded to follow more complicated instructions. For example, a mainstream approach to generating high-quality question-answer pair data is still prompting closed-source GPT-4V because of its advanced instruction-following capabilities, while other models generally fail to achieve.

\item There is still a large space for improvement in techniques like M-ICL and M-CoT. Current research on the two techniques is still rudimentary, and the related capabilities of MLLMs are weak. Thus, explorations of the underlying mechanisms and potential improvement are promising.

\item Developing embodied agents based on MLLMs is a heated topic. It would be meaningful to develop such agents that can interact with the real world. Such endeavors require models with critical capabilities, including perception, reasoning, planning, and execution. 

\item Safety issues. Similar to LLMs, MLLMs can be vulnerable to crafted attacks~\cite{jeong2023hijacking,zhao2023evaluating,shayegani2023jailbreak}. In other words, MLLMs can be misled to output biased or undesirable responses. Thus, improving model safety will be an important topic.

\end{itemize}
